\documentclass[12pt]{article}
\usepackage[utf8]{inputenc}
\usepackage[T1]{fontenc}
\usepackage{hyperref}
\title{Tutong 1 Language Resources for K-12 Curriculum Development}
\author{Compiled by Claude Sonnet 5 (claude-sonnet-5) for Andr\'as Kornai}
\date{}
\begin{document}
\maketitle

\section{Introduction}
The identification given for this language (ISO 639-3 ttx; Glottocode tuto1242) lists the family as ``Bookkeeping''. That is a placeholder rather than a language family, so no family is assumed here. A search turned up a Brunei documentation project (2007--2009, funded by UBD) stating that the Ethnologue label ``Tutong 2'' is misleading. It also says that Wurm and Hattori's ``Tutong 1'' is called the Dusun language in Brunei and is never called ``Tutong'' there. The Tutong language proper (ISO 639-3 ttg) is a different language, and the resources below should not be treated as belonging to ttx. I could not confirm the identity of ttx beyond this one search summary.

\section{Primary Educational Resources}
None were found.

\section{Linguistic and Reference Materials}
\begin{itemize}
\item Tutong language (Wikipedia). This is about ttg, a different language, and is listed only for disambiguation: \url{https://en.wikipedia.org/wiki/Tutong_language}
\item Languages of Brunei (Wikipedia): \url{https://en.wikipedia.org/wiki/Languages_of_Brunei}
\end{itemize}
The search summary reports that Brunei's language agency, Dewan Bahasa dan Pustaka, published a Tutong--Malay dictionary in 1991 and a multilingual word list in 2011. These concern the Tutong language, and I saw no URLs for them.

\section{Religious and Community Resources}
None were found.

\section{Digital and Audio Resources}
\begin{itemize}
\item PARADISEC item from the Brunei language documentation recordings (BN1/019): \url{https://catalog.paradisec.org.au/repository/BN1/019}
\end{itemize}

\section{Curriculum Development Context}
The identity of this language needs checking first, in particular whether it is the Dusun language of Brunei. Nothing found shows any school use. Little is available, and any resources would have to come from the Brunei documentation work and national language agencies.
\end{document}